\section{Natural Language Sentences}
%what did we discover during the interviews? what characteristic should the system have?
%Notice that, the description here, although rich, is at high level.

To bridge the gap between students and teachers in teaching activities, supply a modern efficient way to achieve higher quality study aim, our group is ready to build an advanced, online teaching information system to help students learn well and teachers teach better.\\


From technology background, this information system is based on MySQL database with a front-end page in Vue or Angular or React, while in its back-end side we use Java and Spring frameworks. Running in school local computer room or cloud web servers like AWS, Azure etc.\\


As its functions our group ready to explain with 3 roles, that is students, teachers, and operation officers, these three parts make up together a whole part.\\


For students, first students need to log in with school account or special system account. After log in, there will be a main home page to show all the lectures that teaching in this school. Students need to select their interesting courses to make the course under their plan. Teachers can set enrolment key for their private lectures to keep its resources safe. After students add the lectures, there will be a new page with the title of lecture and teaching teacher’s name which can be arrange with the teacher. Different document type like pdf file, http link and video etc. can be arrange in this page. Besides the course resources, a row of other function can be listed in the course page like forum, multimedia resources, participants etc. \\


As for its teacher side, it can be adjusted to generate the course main page by the teacher. There are also many rich components can be use in the teachers. Just like attendance function, teacher can ask students to scan QR code or tap a key to make sure how many students had attended this course. The teacher side could clearly show the quantity of real attendance students. And one teacher could have many lectures, like one teacher could have machine learning and computer vison these two lecture, he can choose the course he would like to edit. Also, for future perspective, we can make a quiz sub-system for the main system, it can be used inside the main page, a course could have slides, video and chapter quiz to check the quality of teaching.\\


To its operation offices side, administrator can check the full list of all school course, it has a view screen to show the number of courses, number of students, all the videos uploaded to system. It should have an overall view to the whole database of this system.\\



\begin{itemize}
    \item \textbf{Student login page}: can merge with the teacher login page, both solution all could be adopt. 
    \item \textbf{Student main page}: list with participant, video gallery, grades
    \item \textbf{Teacher login page}: as above
    \item \textbf{Teacher main page}: different lectures
    \item \textbf{Operation officers page}: an overall view of the whole system's database
   
 
\end{itemize}
